\documentclass[10pt]{article}

\usepackage[utf8]{inputenc}

\usepackage[T1]{fontenc}

\usepackage{amsmath}

\usepackage{amsfonts}

\usepackage{amssymb}

\usepackage[version=4]{mhchem}

\usepackage{stmaryrd}

\usepackage{bbold}

\usepackage{graphicx}

\usepackage[export]{adjustbox}

\graphicspath{ {./images/} }



\begin{document}

Здесь $E[\Phi]-$ среднее от нек-рой случайной величины $\Phi$, заданной на нек-ром вероятностном пространстве. Было выяснено, что величина Ф(.) должна быть евклидово-ковариантным марковским случайным полем (см. Марковский процесс), удовлетворяющим определенным дополнительным требованиям, и доказано, что эти требования эквивалентны условиям Остервальдера-Шрадера. Во-вторых, переход к евклидовым вероятностным мерам позволил при исследовании проблем, связанных со снятием обрезаний, использовать корректные интегральные представления для регуляризованных функщий Швингера $S_{n}(\cdot \mid \Lambda, \sigma)$ и многочисленные методы статистич. физики. Тем самым оказалось, что многие вопросы евклидовой теории эквивалентны проблемам статистич. физики. На этом пути были упрощены доказательства существования функций Вайтмана в модели $P^{n}(\varphi)_{2}$ и доказано их сушествование в модели $\varphi_{3}^{4}$.



Математически развитые методы К. К. Т. П. привели пока к доказательству существования квантованных гейзенберговых полей только в суперперенормируемых моделях квантовой теории поля. Переход к изучению перенормируемых моделей требует привлечения новых идей и методов.



Лит.: [1] С а й м о н Б., Модель $P(\varphi)_{2}$ евклидовой квантовой теории поля, пер. с англ., М., 1976; [2] Конструктивная теория поля, пер. с англ., М., 1977; [3] Евклидова квантовая теория поля. Марковский подход, М., 1978; [4 ] Гли м м Д., Джаффе А., Математические методы квантовой физики. Подход с использованием функциональных интегралов, пер. с англ., М., $1984 . \quad$ Г. В. Ефимов. КОНТАКТИЗАЦИЯ - см. Контактная геометрия, Симплектическая геометрия.



КОНТАКTHАЯ ГEOMETPИЯ - нечетномерный аналог симплектической геометрии. Взаимоотношения между ними аналогичны взаимоотношениям между геометрией проективного и векторного пространства.



Основные определения. К о н т а К т н о й с т р у к т у р о й в нечетномерном многообразии $M^{2 n+1}$ называется невырожденное $2 n$-мерное распределение, то есть такое распределение, к-рое локально можно задать как ядро 1-формы $\lambda$, удовлетворяющей условию невырожденности: $\lambda \wedge(d \lambda)^{n} \neq 0$ (ни в одной точке).



Если на многообразии задана контактная структура, то многообразие называется контактным многообра 3 и е м; 1-форма, задающая контактную структуру, называется контакт ной фор м ой (такая всюду определенная форма существует не для всякой контактной структуры). Контактные формы, задающие одну и ту же контактную структуру, отличаются умножением на функцию, нигде не обращающуюся в нуль. Если $\lambda-$ контактная форма, то 2 -форма $d \lambda$ определяет на $2 n$-мерной контактной гиперплоскости невырожденную 2-форму, то есть задает структуру симплектич. пространства. Это условие в свою очередь эквивалентно тому, что распределение, заданное 1-формой $\lambda$, невырожденно, то есть является контактной структурой.



Основные примеры. Первый пример контактной структуры - это двумерное распределение в $\mathbb{R}^{3}$, заданное контактной формой $\lambda=x d y+d z(x, y$ и $z-$ декартовы координаты). Подобная структура существует и в $\mathbb{R}^{2 n+1}$ :



\[

\lambda=\sum x_{i} d y_{i}+d z

\]



где $x_{1}, \ldots, x_{n}, y_{1}, \ldots, y_{n}, z$ - декартовы координаты. Локально все контактные структуры в $(2 n+1)$-мерном многообразии диффеоморфны этой (контактная теорема Дарбу).



Стандартная контактная структура в трехмерной сфере $S^{3}$ задается, как распределение комплексных прямых $\mathbb{C}^{1}$, касательных к сфере $S^{3}$, вложенной в $\mathbb{C}^{2}$ в качестве единичной сферы. Эквивалентное описание: контактное распределение состоит из плоскостей, перпендикулярных слоям расслоения Хопфа (то есть окружностям, являюшимся орбитами действия группы $U(1)$ на $S^{3} \subset \mathbb{C}^{2}$ ).



284 КОНСТРУКТИВНАЯ Аналогично строится стандартная контактная структура в любой нечетномерной сфере.



Важным примером контактного многообразия служит многообразие 1-струй функций $J M$ многообразия $M^{n}$, состоящее из 1-струй



\[

(x, f(x), d f(x)), x \in M .

\]



Многообразие $J^{1} M=\mathbb{R} \times T^{*} M$ имеет контактную форму $\lambda=d z-\alpha$, где $z-$ координата на оси $\mathbb{R}$ значений функций, а $\alpha=\sum p_{i} d q_{i}-$ форма действия в $T^{*} M$.



Само название «контактная структура» обязано существованием другому примеру - многообразию контактных элементов. Контактным элеме нтом на многообразии $M^{n}$ называется гиперплоскость в касательном пространстве к $M$ в нек-рой точке, называемой точ о й к о н т а к т а. Различают многообразие всех контактных элементов, являющееся проективизацией кокасательного расслоения: $P\left(T^{*} M\right)$, и двулистно накрывающее его многообразие ориентированных контактных элементов, являющееся сферизацией кокасательного расслоения.



Многообразие контактных элементов несет контактную структуру. Именно, касательный вектор к этому многообразию в точке $\xi$, где $\xi-$ контактный элемент, лежит в контактной гиперплоскости, если его проекция в TM лежит в гиперплоскости $\xi$. Более наглядно: перемещение контактного элемента касается контактной гиперплоскости, если скорость точки контакта принадлежит этому контактному элементу.



Симллектизация и контактизация. Связь между контактной и симплектич. геометрией устанавливается двумя конструкциями: симплектизацией и контактизацией.



Если $M-$ контактное многообразие, то над $M$ существует одномерное расслоение, слой к-рого над точкой $x \in M$ состоит из ненулевых функционалов на $T_{x} M$, равных нулю на контактной гиперплоскости. На пространстве $N$ этого расслоения определена 1-форма $\alpha$; еe значение на векторе $\boldsymbol{v} \in T M$ в точке $\xi \in N$ равно значению функционала $\xi$ на проекции $v$ в $T M$. 2-форма $d \alpha$ определяет в $N$ симплектич. структуру; многообразие $N$ называется с и мпл е к т изац и е й контактного многообразия $M$.



Пример симплектизации: если в качестве $M$ взять многообразие контактных элементов $P T^{*} X$, то симплектизацией $N$ будет $T^{*} X-X-$ кокасательное расслоение без нулевого сечения. На многообразии $N$ действует мультипликативная группа $\mathbb{R}^{x}$ ненулевых действительных чисел умножением функционалов на скаляры. Симплектич. структура $d \alpha$ в $N$ однородна степени 1 относительно этого действия.



Симплектизация позволяет сводить вопросы, относящиеся к К. г., к аналогичным вопросам симплектич. геометрии. Напр., контактоморфизмы многообразия $M$ (то есть диффеоморфизмы, сохраняющие контактную структуру) взаимно однозначно соответствуют симплектоморфизмам многообразия $N$, коммутируюшим с действием группы $\mathbb{R}^{x}$.



Контактизация сопоставляет точному симплектич. многообразию $N$ (то есть многообразию, на к-ром симплек-



\includegraphics[max width=\textwidth, center]{2023_07_08_94edad980172c6403d7fg-1(1)}

контактное многообразие $K=N \times \mathbb{R}$ с контактной формой $\alpha-d z$, где $z-$ координата в $\mathbb{R}$.



Напр., если $N$ является кокасательным расслоением $T^{*} M$,



\includegraphics[max width=\textwidth, center]{2023_07_08_94edad980172c6403d7fg-1}

Контактизация определена не однозначно - она зависит от выбора 1-формы $\alpha$ в равенстве $\omega=d \alpha$.



Контактные векторные поля. К о н т а т н о е вект о р но е поле - это такое векторное поле на контактном многообразии, локальный поток к-рого сохраняет контактную структуру, то есть является контактоморфизмом. Контактные векторные поля образуют алгебру Ли.



Пример такого поля - геодезич. поток. Это векторное поле определено в многообразии контактных элементов риманова многообразия. Его поток параллельно переносит контактные элементы вдоль геодезических, выходящих из





\end{document}